\begin{tcolorbox}[enhanced,colback=paper,colframe=linegray,
  boxrule=.7pt,arc=2mm,left=3mm,right=3mm,top=3mm,bottom=2.5mm,
  fonttitle=\sffamily\bfseries,coltitle=ink,
  title={图解：RNN 将同一个更新函数沿时间展开}]
\begin{center}
\begin{tikzpicture}[>=Latex,
  state/.style={draw=blue,fill=blueback,rounded corners=2.5pt,
    minimum width=23mm,minimum height=12mm,align=center,font=\small},
  word/.style={draw=linegray,fill=white,rounded corners=2.5pt,
    minimum width=23mm,minimum height=9mm,align=center,font=\small}]
\node[word,minimum width=18mm] (h0) at (-1.4,0) {$h_0=0$};
\node[state] (h1) at (1.75,0) {更新函数\\$h_1$};
\node[state] (h2) at (5.65,0) {同一函数\\$h_2$};
\node[state] (h3) at (9.55,0) {同一函数\\$h_3$};
\node[word] (x1) at (1.75,-1.55) {春天\\$x_1$};
\node[word] (x2) at (5.65,-1.55) {花\\$x_2$};
\node[word] (x3) at (9.55,-1.55) {开了\\$x_3$};
\foreach \a/\b in {x1/h1,x2/h2,x3/h3}
  \draw[->,very thick,blue] (\a)--(\b);
\draw[->,very thick,orange] (h0)--node[above,font=\scriptsize]{旧状态}(h1);
\draw[->,very thick,orange] (h1)--node[above,font=\scriptsize]{传递状态}(h2);
\draw[->,very thick,orange] (h2)--node[above,font=\scriptsize]{传递状态}(h3);
\end{tikzpicture}
\end{center}
\small\noindent
\textbf{\color{blue}蓝色：}当前 Token 进入更新函数；
\textbf{\color{orange}橙色：}上一时刻的状态传给下一时刻。
三个蓝框\textbf{共享一组可学习参数}，并非训练了三个不同的网络。
第 $3$ 步必须等到 $h_2$ 算出来，才能继续计算 $h_3$。
\end{tcolorbox}
